% Capstone 6: elliptic-curve arithmetic. Port of «Мини-курсовая - арифметика эллиптических кривых».
% The performance numbers come from the generated requirements.tex (make bench).
\documentclass{ca-capstone}
\capstonerefs{sagemath,gmp,gmpy2,crandall-pomerance,hps}
\begin{document}
% Generated by `make bench CAPSTONE=ec-arithmetic` (taskgen.capstone report) from
% bench/results.json: reference time x slack on the reference machine.
% Do not edit: change bench/config.yaml and re-run the benchmark.
\newcommand{\ReqSlack}{10}
\newcommand{\ReqBitsMin}{512}
\newcommand{\ReqBitsMax}{1024}
\newcommand{\ReqAddUs}{200}
\newcommand{\ReqMulMs}{500}
\newcommand{\ReqPointMs}{1000}
\newcommand{\ReqOrderBits}{64}
\newcommand{\ReqOrderSeconds}{5}

\capstonetitle{Арифметика эллиптических кривых}

\goal
Требуется реализовать библиотеку для работы в группе точек на эллиптической кривой по модулю
простого числа $\EC(\GF{p})$.

\functionality
\begin{enumerate}
  \item Сложение точек, нахождение обратной точки
  \item Умножение точки на число (<<возведение в степень>>) с помощью алгоритма удваивания точки
  \item Умножение точки на число (<<возведение в степень>>) на основе тернарного разложения
    показателя
  \item Порождение точки на заданной кривой\footnote{т.\,е. нахождение такой точки
    $P = (x_p, y_p)$, что $x_p, y_p$ удовлетворяют уравнению заданной кривой}
  \item Подсчёт порядка (количества точек) кривой методом Шенкса"--~Местре\footnote{см
    п.~7.5.1 в~\cite{crandall-pomerance}}
\end{enumerate}

\requirements
Очевидными требованиями являются:
\begin{enumerate}
  \item корректность результата выполняемых операций
  \item достаточно высокая скорость работы (см. комментарий ниже)
\end{enumerate}

\textbf{Корректность работы} можно проверять, сравнивая результаты работы своих функций
с~результатами функций некоторой библиотеки, например, \emph{SageMath}~\cite{sagemath}.

Приемлемой будет считаться \textbf{скорость работы}, позволяющая выполнять базовые
арифметические операции с эллиптическими кривыми, построенными на поле вычетов по модулю~$N$
длиной \textbf{\ReqBitsMin--\ReqBitsMax~бит}. Ориентиры для модуля длиной \ReqBitsMax~бит:
сложение точек "--- не более \ReqAddUs~мкс, умножение точки на число той же длины "--- не более
\ReqMulMs~мс, порождение точки "--- не более \ReqPointMs~мс; подсчёт порядка кривой методом
Шенкса"--~Местре для модуля длиной \ReqOrderBits~бит "--- не более \ReqOrderSeconds~с.

\details
В работе требуется реализовать базовые операции арифметики эллиптических кривых <<с нуля>>,
поэтому, очевидно, запрещено пользоваться соответствующими функциями специализированных
библиотек. Однако \textbf{использование библиотек} для работы \textbf{с большими числами
разрешено}\footnote{рекомендуется использовать библиотеку \emph{GNU GMP}~\cite{gmp} или её
обёртку~\cite{gmpy2}}.

\textbf{Язык может быть любым}. Например, если ваша реализация на \emph{php} или \emph{JS}
достигает требуемого уровня производительности, выбор языка будет обоснован.

Платформо-независимый, \textbf{переносимый код} писать \textbf{не требуется}, но во время сдачи
должна быть возможность проверить работу библиотеки.

\schedule
Работа выполняется \textbf{индивидуально}. Срок выполнения "--- \textbf{месяц}.

Работа должна быть ваша: реализация должна быть \textbf{разработана и написана вами}, а не
собрана из кусков чужих работ.

Желательно хранить код в \emph{git}-репозитории на \emph{github}, \emph{bitbucket} или подобных
платформах. \textbf{Документация не требуется}.

\remarks
\begin{enumerate}
  \item Нахождение тернарного разложения показателя степени~$n$ требуется реализовать
    самостоятельно, в виде отдельной функции
  \item Для нахождения точки на кривой можно реализовать любой алгоритм, успешно завершающийся
    за приемлемое время, например, подбор координаты~$y$ на основе метода Тонелли"--~Шенкса
  \item Для подсчёта порядка кривой требуется реализовать метод Шенкса"--~Местре, однако
    реализация метода Шуфа (см. п.~7.5.2 в~\cite{crandall-pomerance}) будет существенным плюсом
    вашей работы
  \item Проверку корректности работы реализуемых операций можно осуществлять с помощью т.\,н.
    \emph{Unit-тестов}
\end{enumerate}

\references
\end{document}
